% ============================================================
% Cheat sheet, argomento: il pattern Strategy (dispensa, capitolo 15)
% ============================================================
\argomento{Strategy}{Dispensa: capitolo 15. Pattern comportamentale della GoF}[Pattern Strategy]

\scheda{Carta d'identità}{
  \pintento{Definire una \textbf{famiglia di algoritmi}, incapsularli uno per classe e renderli
    \textbf{intercambiabili}: l'algoritmo può variare indipendentemente da chi lo usa.}
  \ptesto{«modi diversi di calcolare (spedire, pagare, ordinare, scontare, validare)», «politiche», «il cliente
    sceglie il metodo di spedizione», «cambiare l'algoritmo a tempo di esecuzione», «aggiungere un nuovo modo di
    calcolare senza modificare la classe che lo usa».}
  \pcodice{un'interfaccia con \textbf{un} metodo (o pochi), che è l'algoritmo (\code{double costo(double pesoKg)}); più
    classi che la implementano, ognuna con la sua formula; il \textbf{contesto} ha un campo di quel tipo,
    \textbf{impostato dal client} (costruttore o \code{setSpedizione(\ldots)}), e gli delega il calcolo. Spesso al posto
    delle classi c'è una \textit{lambda} o un \code{Comparator}.}
  \pdiagramma{\code{Context} con un'aggregazione verso \code{«interface» Strategy}, realizzata da
    \code{ConcreteStrategyA}, \code{B}, \code{C}: la stessa forma dello State.}
  \pvarianti{la strategia nel costruttore (fissa) o in un setter (cambia quando il client vuole); con una lambda
    (\code{ordine.setSpedizione(peso -> 0.0)}); nella libreria di Java: \code{lista.sort(comparatore)}, dove il
    \code{Comparator} è la strategia di confronto.}
  \pconfondere{lo \textbf{State} (\vv{state-struttura}): stessa struttura, ma lì l'oggetto cambia \textbf{da solo},
    perché gli stati chiamano \code{setState}; la \textbf{Dependency Injection} (\vv{di-struttura}): dice \textbf{come}
    arriva l'oggetto (da fuori), non che cosa rappresenta; una strategia di solito arriva proprio per iniezione.}
}

\begin{struttura}{strategy-struttura}{La struttura}
\begin{tikzpicture}
  \umlclass[2.6cm]{c}{(0,0)}{Context}{}{+ esegui()}
  \umlclass[3.0cm]{s}{(4.05,0)}{\stereotipo{interface}{Strategy}}{}{+ algoritmo()}
  \node[classe1=2.7cm] (a) at (1.1,-2.2) {ConcreteStrategyA};
  \node[classe1=2.7cm] (b) at (4.05,-2.2) {ConcreteStrategyB};
  \node[classe1=2.7cm] (cc) at (7.0,-2.2) {ConcreteStrategyC};
  \draw[aggr] (c.east) -- (s.west);
  \estO{(s.west)}{1}{strategia}
  \foreach \x in {a, b, cc}{ \draw[realizza] (\x.north) |- ($(s.south)+(0,-0.45)$) -- (s.south); }
  \node[notauml, anchor=north] (n) at (0,-1.3) {strategia.algoritmo()};
  \draw[tratt] (c.south) -- (n.north);
\end{tikzpicture}
\tcblower
\textbf{Strategy}: l'interfaccia comune degli algoritmi. \textbf{ConcreteStrategy}: un algoritmo ciascuna; non si
conoscono fra loro. \textbf{Context}: usa una strategia senza sapere quale; la sceglie il \textbf{client}.

\renewcommand{\arraystretch}{1.15}\footnotesize
\begin{tabular}{@{}l l l@{}}
   & \textbf{State} & \textbf{Strategy} \\ \hline
  chi cambia l'oggetto & gli stati & il client \\
  si conoscono? & sì & no \\
  rappresenta & una fase & un modo di fare \\
\end{tabular}
\end{struttura}

\begin{programma}{strategy-esempio}{L'esempio della dispensa: il costo della spedizione}
\begin{lstlisting}[style=javacs, name=Negozio.java]
public class Negozio {
    public static void main(String[] args) {
        Ordine ordine = new Ordine(40.0, 6.0);
        System.out.println("Standard:   "
                + ordine.totale());
        ordine.setSpedizione(new Express());
        System.out.println("Express:    "
                + ordine.totale());
        ordine.setSpedizione(new RitiroInNegozio());
        System.out.println("In negozio: "
                + ordine.totale());
    }
}
interface StrategiaSpedizione {         // Strategy
    double costo(double pesoKg);
}
class Standard implements StrategiaSpedizione {
    public double costo(double pesoKg) {
        return 5.0 + 0.5 * pesoKg;
    }
}
class Express implements StrategiaSpedizione {
    public double costo(double pesoKg) {
        return 12.0 + 1.0 * pesoKg;
    }
}
class RitiroInNegozio
        implements StrategiaSpedizione {
    public double costo(double pesoKg) {
        return 0.0;
    }
}
class Ordine {                          // Context
    private final double merce, pesoKg;
    private StrategiaSpedizione spedizione =
            new Standard();
    Ordine(double merce, double pesoKg) {
        this.merce = merce;
        this.pesoKg = pesoKg;
    }
    void setSpedizione(StrategiaSpedizione s) {
        spedizione = s;      // la sceglie il client
    }
    double totale() {
        return merce + spedizione.costo(pesoKg);
    }
}
\end{lstlisting}
\tcblower
\begin{tikzpicture}
  \umlclass[3.0cm]{o}{(0,0)}{Ordine}{- merce: double\\- pesoKg: double}{\textasciitilde\ setSpedizione(s)\\\textasciitilde\ totale(): double}
  \umlclass[3.2cm]{s}{(4.1,0)}{\stereotipo{interface}{StrategiaSpedizione}}{}{+ costo(pesoKg): double}
  \node[classe1=1.6cm] (a) at (2.0,-2.6) {Standard};
  \node[classe1=1.5cm] (b) at (3.75,-2.6) {Express};
  \node[classe1=2.3cm] (c) at (5.9,-2.6) {RitiroInNegozio};
  \draw[aggr] (o.east) -- (o.east -| s.west);
  \estO{(o.east -| s.west)}{1}{spedizione}
  \foreach \x in {a, b, c}{ \draw[realizza] (\x.north) |- ($(s.south)+(0,-0.45)$) -- (s.south); }
\end{tikzpicture}

\medskip
\begin{lstlisting}[style=shellcs, name=Terminale]
$ javac Negozio.java
$ java Negozio
Standard:   48.0
Express:    58.0
In negozio: 40.0
\end{lstlisting}
\regola{Lo stesso ordine (40 euro di merce, 6 kg) cambia algoritmo tre volte, e \code{Ordine} non cambia mai: per una
nuova spedizione basta una classe nuova. Le slide nominano lo Strategy nella tabella dei pattern senza un esempio;
questo è quello della dispensa.}
\end{programma}
